% Preámbulo compartido para todas las clases de Beamer.
% Cada clase lo incluye con % Preámbulo compartido para todas las clases de Beamer.
% Cada clase lo incluye con % Preámbulo compartido para todas las clases de Beamer.
% Cada clase lo incluye con \input{../../latex/preamble.tex}

\documentclass[aspectratio=169,11pt]{beamer}

\usetheme{Madrid}
\setbeamertemplate{navigation symbols}{}
\setbeamertemplate{footline}[frame number]

\usepackage[utf8]{inputenc}
\usepackage[spanish]{babel}
\usepackage{amsmath,amssymb}
\usepackage{graphicx}
\usepackage{booktabs}
\usepackage{listings}
\usepackage{xcolor}

\lstset{
  basicstyle=\ttfamily\small,
  keywordstyle=\color{blue},
  commentstyle=\color{gray},
  stringstyle=\color{teal},
  showstringspaces=false,
  breaklines=true,
  frame=single
}

\newcommand{\definicion}[1]{%
  \begin{block}{Definición}
    #1
  \end{block}
}

\newcommand{\ejemplo}[1]{%
  \begin{exampleblock}{Ejemplo}
    #1
  \end{exampleblock}
}

\newcommand{\importante}[1]{%
  \begin{alertblock}{Importante}
    #1
  \end{alertblock}
}

\author{Agustín Gurvich}
\institute{Introducción a Machine Learning}


\documentclass[aspectratio=169,11pt]{beamer}

\usetheme{Madrid}
\setbeamertemplate{navigation symbols}{}
\setbeamertemplate{footline}[frame number]

\usepackage[utf8]{inputenc}
\usepackage[spanish]{babel}
\usepackage{amsmath,amssymb}
\usepackage{graphicx}
\usepackage{booktabs}
\usepackage{listings}
\usepackage{xcolor}

\lstset{
  basicstyle=\ttfamily\small,
  keywordstyle=\color{blue},
  commentstyle=\color{gray},
  stringstyle=\color{teal},
  showstringspaces=false,
  breaklines=true,
  frame=single
}

\newcommand{\definicion}[1]{%
  \begin{block}{Definición}
    #1
  \end{block}
}

\newcommand{\ejemplo}[1]{%
  \begin{exampleblock}{Ejemplo}
    #1
  \end{exampleblock}
}

\newcommand{\importante}[1]{%
  \begin{alertblock}{Importante}
    #1
  \end{alertblock}
}

\author{Agustín Gurvich}
\institute{Introducción a Machine Learning}


\documentclass[aspectratio=169,11pt]{beamer}

\usetheme{Madrid}
\setbeamertemplate{navigation symbols}{}
\setbeamertemplate{footline}[frame number]

\usepackage[utf8]{inputenc}
\usepackage[spanish]{babel}
\usepackage{amsmath,amssymb}
\usepackage{graphicx}
\usepackage{booktabs}
\usepackage{listings}
\usepackage{xcolor}

\lstset{
  basicstyle=\ttfamily\small,
  keywordstyle=\color{blue},
  commentstyle=\color{gray},
  stringstyle=\color{teal},
  showstringspaces=false,
  breaklines=true,
  frame=single
}

\newcommand{\definicion}[1]{%
  \begin{block}{Definición}
    #1
  \end{block}
}

\newcommand{\ejemplo}[1]{%
  \begin{exampleblock}{Ejemplo}
    #1
  \end{exampleblock}
}

\newcommand{\importante}[1]{%
  \begin{alertblock}{Importante}
    #1
  \end{alertblock}
}

\author{Agustín Gurvich}
\institute{Introducción a Machine Learning}


\title{Introducción a ML \\ y Probabilidad/Estadística}
\subtitle{¿Qué es aprender de datos? Matemática básica}
\date{}

\begin{document}

\frame{\titlepage}

\part{Introducción a Machine Learning}
\frame{\partpage}

\section{¿Qué es Machine Learning?}

\begin{frame}{¿Qué es Machine Learning?}
  \definicion{
    El machine learning es un subconjunto de la inteligencia artificial centrado en algoritmos que pueden \textbf{aprender} los patrones de los datos de entrenamiento y hacer inferencias precisas sobre nuevos datos.
  }
  \pause
  En vez de escribir reglas a mano, le mostramos ejemplos al modelo
  y este \textbf{aprende un patrón} que generaliza a casos nuevos.
\end{frame}

\begin{frame}{Programación tradicional vs. ML}
  \begin{columns}
    \column{0.5\textwidth}
    \textbf{Programación tradicional}
    \begin{itemize}
      \item Le damos datos y definimos reglas. Eso nos da \textit{resultados}
      \item Nosotros escribimos la lógica
    \end{itemize}
    \column{0.5\textwidth}
    \textbf{Machine Learning}
    \begin{itemize}
      \item Le damos datos y resultados. Eso nos da \textit{reglas}
      \item El modelo aprende la lógica
    \end{itemize}
  \end{columns}
\end{frame}

\begin{frame}{¿Dónde se usa?}
  \begin{itemize}
    \item Recomendaciones (Netflix, Spotify, YouTube)
    \item Detección de spam
    \item Reconocimiento de imágenes y voz
    \item Diagnóstico médico asistido
    \item Autos autónomos
    \item Traducción automática
  \end{itemize}
\end{frame}

\section{Tipos de aprendizaje}

\begin{frame}{Aprendizaje supervisado}
  \definicion{
    Aprendemos una función que asocia entradas a salidas, a partir de
    ejemplos \textbf{etiquetados}.
  }
  \ejemplo{
    Predecir el precio de una casa a partir de sus características
    (m², ubicación, antigüedad), usando casas ya vendidas como ejemplo.
  }
  Se suelen usar en dos tipos de problemas: \textbf{clasificación} (categorías) y
  \textbf{regresión} (valores numéricos continuos).
\end{frame}

\begin{frame}{Aprendizaje no supervisado}
  \definicion{
    Buscamos patrones o estructura en datos \textbf{sin etiquetas}.
  }
  \ejemplo{
    Agrupar clientes de una tienda según sus hábitos de compra, sin
    saber de antemano qué grupos existen (clustering).
  }
\end{frame}

\begin{frame}{Aprendizaje por refuerzo}
  \definicion{
    Un agente aprende a tomar decisiones interactuando con un entorno,
    recibiendo recompensas o castigos.
  }
  \ejemplo{
    Un programa que aprende a jugar al go o a un videojuego
    jugando muchas partidas contra sí mismo.
  }
\end{frame}

\section{Flujo de trabajo en ML}

\begin{frame}{¿Cómo es el flujo de trabajo en un proyecto de ML?}
  ¿Cómo se imaginan que trabaja un ingeniero de ML?
  \pause
  \begin{enumerate}
    \item Definir el problema
    \item Recolectar y explorar los datos
    \item Preprocesar los datos (limpieza, transformación)
    \item Elegir un modelo y entrenarlo
    \item Evaluar el modelo
    \item Ajustar (iterar) y poner en producción
  \end{enumerate}
  \pause
  \importante{
    La mayor parte del tiempo real de un proyecto de ML se va en los
    datos, no en el modelo.
  }
\end{frame}

\begin{frame}{Palabras que van a escuchar mucho}
  \begin{itemize}
    \item \textbf{Dataset}: conjunto de datos con los que trabajamos
    \item \textbf{Feature}: cada variable de entrada
    \item \textbf{Target (etiqueta)}: lo que queremos predecir
    \item \textbf{Modelo}: la función/algoritmo que aprende el patrón
    \item \textbf{Entrenamiento}: proceso de ajustar el modelo a los datos
    \item \textbf{Overfitting}: el modelo memoriza en vez de generalizar
  \end{itemize}
\end{frame}

\part{Probabilidad y Estadística para ML}
\frame{\partpage}

\section{¿Por qué estadística en ML?}

\begin{frame}{¿Por qué necesitamos estadística?}
  \importante{
    ML (y la computación en general) es, en el fondo, matemática aplicada a gran escala. Todo dataset
    es una muestra, y todo modelo hace operaciones estadísticas sobre esa muestra.
  }
  \begin{itemize}
    \item Para \textbf{entender} los datos antes de modelar (¿hay outliers? ¿están sesgados?)
    \item Para \textbf{elegir} qué transformar y cómo
    \item Para \textbf{interpretar} si un resultado es significativo o casualidad
  \end{itemize}
\end{frame}

\section{Medidas de tendencia central y dispersión}

\begin{frame}{Medidas de tendencia central}
  \begin{itemize}
    \item \textbf{Media ($\bar{x}$)}: promedio de los valores
      $$\bar{x} = \frac{1}{n}\sum_{i=1}^{n} x_i$$
    \item \textbf{Mediana}: valor que queda en el medio al ordenar los datos
    \item \textbf{Moda}: valor más frecuente
  \end{itemize}
  \pause
  \ejemplo{
    Sueldos: 1, 1, 2, 2, 2, 100 (en cientos de miles).
    Media $\approx 18$, mediana $= 2$, moda $= 2$.
    La media se deja engañar por valores extremos, la mediana no.
  }
\end{frame}

\begin{frame}{Medidas de dispersión}
  \begin{itemize}
    \item \textbf{Varianza}: qué tan lejos están los datos de la media, en promedio (al cuadrado)
      $$\sigma^2 = \frac{1}{n}\sum_{i=1}^{n} (x_i - \bar{x})^2$$
    \item \textbf{Desvío estándar}: $\sigma = \sqrt{\sigma^2}$, en las mismas unidades que los datos
  \end{itemize}
  \pause
  \importante{
    Un desvío estándar grande dice que los datos están muy dispersos!
    La media sola no alcanza para describir un dataset.
  }
\end{frame}

\section{Distribuciones y correlación}

\begin{frame}{La distribución normal}
  \begin{itemize}
    \item Forma de \textbf{campana}, simétrica alrededor de la media
    \item Muchísimos fenómenos naturales se aproximan a esta forma (alturas, errores de medición)
    \item Muchos algoritmos de ML \textbf{asumen} (implícita o explícitamente) que las variables se comportan así
  \end{itemize}
\end{frame}

\begin{frame}{Probabilidad: lo mínimo indispensable}
  \begin{itemize}
    \item \textbf{Probabilidad}: número entre 0 y 1 que mide qué tan posible es un evento
    \item \textbf{Probabilidad condicional} $P(A \mid B)$: probabilidad de $A$ sabiendo que $B$ ya ocurrió
  \end{itemize}
  \ejemplo{
    $P(\text{spam})$ = probabilidad general de que un mail sea spam.\\
    $P(\text{spam} \mid \text{contiene ``gratis''})$ = probabilidad de spam
    \emph{sabiendo} que el mail contiene la palabra ``gratis''. Mucho más útil.
  }
\end{frame}

\begin{frame}{Correlación}
  \definicion{
    La correlación mide qué tan relacionadas están dos variables
    numéricas, con un valor entre $-1$ y $1$.
  }
  \begin{itemize}
    \item Cerca de $1$: cuando una sube, la otra también (relación positiva)
    \item Cerca de $-1$: cuando una sube, la otra baja (relación negativa)
    \item Cerca de $0$: no hay relación lineal aparente
  \end{itemize}
  \importante{Correlación no implica causalidad.}
\end{frame}

\begin{frame}{Resumen de la clase}
  \centering
  \Large ¿Preguntas?
\end{frame}

\end{document}
